% Figure gsl: generated by figures/rse_fig_gsl.py; placed where the text first discusses it.
\begin{figure}[!tbp]
  \centering
  \includegraphics[width=\textwidth]{../figures/out/rse_fig_gsl.pdf}
  \caption{\textbf{The same audit at Great Salt Lake, Utah, with a lake-bed control from echo soundings.}
  (a, b)~For the north and south arms, the shoreline displacement $D$ at each level bin (levels on NAVD~88), after each surface's vertical offset fitted to flooded area, against the observation noise floor $Q$ (grey area; the \metric{params.err_pct_hi}th percentile of the distances between the shorelines of two interleaved halves of the acquisitions).
  Coloured lines: the five global grids; open markers on the top edge, levels at which a grid predicts no shoreline in the arm.
  Grey dashed line: the 3DEP control, mostly 2016 lidar.
  Black line: the lake-bed control, built from USGS single-beam echo soundings below the sounded top and from digitised map contours and the National Elevation Dataset above it (offset \metric{rev2.g.lakebed.all.north.shape_offset_m}~m north and \metric{rev2.g.lakebed.all.south.shape_offset_m}~m south); filled where $D > Q$ at the level.
  Shaded columns: the sounded levels, bins \metric{rev2.g.lakebed.sounded.north.shape.bins} in the north arm (\metric{rev2.g.lakebed.coverage.north.sounded.level_low_navd88_m}--\metric{rev2.g.lakebed.coverage.north.sounded.level_high_navd88_m}~m) and \metric{rev2.g.lakebed.sounded.south.shape.bins} in the south arm (\metric{rev2.g.lakebed.coverage.south.sounded.level_low_navd88_m}--\metric{rev2.g.lakebed.coverage.south.sounded.level_high_navd88_m}~m).
  With the offset fitted over all its levels, the lake-bed control exceeds $Q$ at \metric{rev2.g.lakebed.all.north.class.sounded.n_exceeding} of \metric{rev2.g.lakebed.all.north.class.sounded.n_evaluable} sounded levels in the north arm and \metric{rev2.g.lakebed.all.south.class.sounded.n_exceeding} of \metric{rev2.g.lakebed.all.south.class.sounded.n_evaluable} in the south arm.
  (c)~For each arm and surface, the median of $D$ over the evaluable levels (marker) against the median $Q$ over the same levels (grey bar ending in a tick); filled where the surface fails the shoreline test in the arm ($D > Q$ at more than half of the levels).
  The lake-bed control is shown over all its levels (\metric{rev2.g.lakebed.all.north.shape.n_evaluable} north, \metric{rev2.g.lakebed.all.south.shape.n_evaluable} south) and over the sounded levels alone, with its offset refitted on them; on the sounded levels it fails the test in both arms (\metric{rev2.g.lakebed.sounded.north.shape.n_exceeding} of \metric{rev2.g.lakebed.sounded.north.shape.n_evaluable} and \metric{rev2.g.lakebed.sounded.south.shape.n_exceeding} of \metric{rev2.g.lakebed.sounded.south.shape.n_evaluable} levels), so the control itself misplaces the sounded shorelines.
  Logarithmic axes.
  (d)~Heights of each grid over each arm's emerged band (share of the band's cells per half-metre bin of height); dotted line, the gauge level of February 2000 (\metric{rev.g.gsl.north.gauge_feb2000_navd88_m}~m north, \metric{rev.g.gsl.south.gauge_feb2000_navd88_m}~m south), when SRTM flew.}
  \label{fig:rse_gsl}
\end{figure}
